\documentclass[10pt,a4paper]{amsart}
\usepackage[margin=19mm]{geometry}
\usepackage{amsmath,amssymb,mathtools}
\usepackage[scr=rsfs,cal=euler]{mathalfa}
\usepackage{xfrac}
\usepackage[unicode,hidelinks,bookmarks=false]{hyperref}
\newcommand{\ie}{i.e.\ }
\newcommand{\cf}{cf.\ }
\newcommand{\resp}{respectively\ }
\newcommand{\Subsection}{\subsection}
\providecommand{\nobreakdash}{\nobreak\hbox{-}}
\setlength{\emergencystretch}{2em}
\multlinegap0pt
\usepackage{amssymb}
\usepackage{tikz}
\newtheorem{theorem}{Theorem}
\DeclareMathOperator{\NC}{NC}
\DeclareMathOperator{\wt}{wt}
\title{On the structure of the $d$-indivisible noncrossing partition posets}
\author{Richard Ehrenborg, G\'abor Hetyei}
\date{Source: arXiv:2407.08577v3 (CC BY 4.0)}
\begin{document}
\maketitle

\noindent\emph{Source and licence.} On the structure of the $d$-indivisible noncrossing partition posets. Version arXiv:2407.08577v3 (CC BY 4.0), distributed by arXiv under the Creative Commons Attribution 4.0 International licence, http://creativecommons.org/licenses/by/4.0/, as recorded in the arXiv licence metadata for this exact version and shown on the abstract page https://arxiv.org/abs/2407.08577v3. Full licence notice: \texttt{licenses/arxiv-2407.08577-CC-BY-4.0.txt}.\par\smallskip

\noindent\textit{Editorial scope.} Counted unit: the article's theorem equation\_C\_and\_C* (source0.tex lines 815--831). The article's complete original proof of it is delivered with it (source0.tex lines 832--991, one \texttt{proof} environment of 160 lines). No mathematics is written, repaired or re-derived: the statement and the proof are the article's own lines, in the article's own notation, and no retained line is substituted or re-flowed.  Retained uncounted as the source-local context the proof needs: lines 721--733.  Nothing was omitted from any retained span, so the package carries no omission record.  No retained line contains a citation, so the wrapper carries no bibliography and no bibliography entry.  No retained line contains a figure environment, an image-inclusion command or drawing code, so no image is delivered and the package declares no asset dependency.  Editorial formatting only, in addition: an AMS \texttt{amsart} preamble that restates the article's own declarations and macros used by the retained lines, namely the environments \texttt{theorem} on separate counters, together with the standard packages those lines need.  The retained environments print in this document as Theorem 1 in order of appearance, whereas the article prints the same items with the numbering of its own full document.  The complete native source of the article as distributed in the arXiv e-print for this exact version is bundled unchanged as \texttt{sources/162187/source0.tex}.
\begin{theorem}[M\"uhle--Nadeau--Williams]
\label{theorem_characterization}
The noncrossing partition $\pi$ belongs to $\NC^{d}_{n}$
if and only if
the following two conditions are satisfied
(i)
$n \equiv 1 \bmod d$;
(ii)
for each pair of cyclically consecutive elements $i$ and $j$
of a nonsingleton block $B$ in a noncrossing partition $\pi$,
the number of elements in $[n]$ strictly between $i$ and $j$ in the cyclic order
is divisible by~$d$.
\end{theorem}

\begin{theorem}
The following three generating function identities hold:
\begin{align}
\label{equation_C_and_C*}
C(x)
& =
A^{*}\left(x \cdot s \cdot C^{*}(x)^{d}\right) ,
&
C^{*}(x)
& =
A\left(x \cdot t \cdot C(x)^{d}\right) , \\
B(x) & = C(x) \cdot C^{*}(x) .
\end{align}
Furthermore, the two identities in~\eqref{equation_C_and_C*} uniquely determine
the generating functions $C(x)$ and $C^{*}(x)$.
\label{theorem_generating_functions_C_and_C*}
\end{theorem}

\begin{proof}
Observe that the first two identities are dual to each other. Hence it
is enough 
that we prove the second identity.
Let $\pi$ be a partition in $\NC^{d}_{dk+1}$ where $dk+1 = n$
and where the singleton block~$\{n^{\prime}\}$
belongs to the dual partition $\pi^{\prime}$.
That is, the partition contains a block $B$ such that
\begin{align*}
B
& =
\{1 = i_{1} < i_{2} < \cdots < i_{dr+1} = n\} .
\end{align*}
By Theorem~\ref{theorem_characterization}
each difference $i_{j+1} - i_{j}$ is congruent to $1$ modulo $d$.
Let $\tau_{j}$ now be the restriction of the partition $\pi$
to the union $[i_{j},i_{j+1}] \cup B$ and identify all the elements of the block $B$.
This construction is equivalent to letting $\tau_{j}$ be the restriction to the
open interval $(i_{j},i_{j+1})$ and adding a singleton block.
Keeping track of the blocks in~$\pi$ we have that
\begin{align*}
|\pi| - 1
& =
(|\tau_{1}| - 1)
+ \cdots +
(|\tau_{dr}| - 1) .
\end{align*}
Let $\tau_{i}^{\prime}$ be the dual partition of $\tau_{i}$ inside the set $[i_{j},i_{j+1}] \cup B$.
Counting the blocks in the dual partition we obtain
\begin{align*}
|\pi^{\prime}| - 1
& =
|\tau^{\prime}_{1}|
+ \cdots +
|\tau^{\prime}_{dr}|
=
dr
+
(|\tau^{\prime}_{1}| - 1)
+ \cdots +
(|\tau^{\prime}_{dr}| - 1) .
\end{align*}
Hence the weight $\wt(\pi)$ factors as
\begin{align*}
\wt(\pi)
& =
a(r)
\cdot
t^{r}
\cdot
\wt(\tau_{1})
\cdots
\wt(\tau_{dr}) .
\end{align*}
Now we express $c^{*}(k)$ as the sum
\begin{align}
\label{equation_recursion_for_c*}
c^{*}(k)
& =
\sum_{r \geq 1}
a(r)
\cdot
t^{r}
\cdot
\sum_{
\substack{1 = i_{1} < i_{2} < \cdots < i_{dr+1} = n \\
i_{j+1} - i_{j} \equiv 1 \bmod d}
}
\sum_{\tau_{1}}
\cdots
\sum_{\tau_{dr}}
\prod_{j=1}^{dr}
\wt(\tau_{j}) \\
\nonumber
& =
\sum_{r \geq 1}
a(r)
\cdot
t^{r}
\cdot
\sum_{
\substack{1 = i_{1} < i_{2} < \cdots < i_{dr+1} = n \\
i_{j+1} - i_{j} \equiv 1 \bmod d}
}
\prod_{j=1}^{dr}
c\left( \frac{i_{j+1} - i_{j} - 1}{d} \right) ,
\end{align}
where each $\tau_{j}$ ranges over partitions in 
$\NC^{d}_{i_{j+1}-i_{j}}$
with a given singleton block.
Observe that $k$ is given by the sum
\begin{align*}
k
& =
r + \frac{i_{2} - i_{1} - 1}{d}
+ \cdots + \frac{i_{dr+1} - i_{dr} - 1}{d} .
\end{align*}
Hence multiply with $x^{k}$ and sum over all $k \geq 1$.
\begin{align*}
\sum_{k \geq 1} c^{*}(k) \cdot x^{k}
& =
\sum_{r \geq 1}
a(r) \cdot t^{r} \cdot x^{r} \cdot
\sum_{
\substack{1 = i_{1} < i_{2} < \cdots < i_{dr+1} = n \\
i_{j+1} - i_{j} \equiv 1 \bmod d}
}
\prod_{j=1}^{dr}
c\left( \frac{i_{j+1} - i_{j} - 1}{d} \right) 
\cdot
x^{({i_{j+1} - i_{j} - 1})/{d}} \\
& =
\sum_{r \geq 1}
a(r) \cdot t^{r} \cdot x^{r} \cdot C(x)^{dr} .
\end{align*}
Lastly, adding the constant $c^{*}(0) = 1$ yields the second equation.

Equation~\eqref{equation_recursion_for_c*} and its dual
yield a pair of intertwined recursions for the sequences
$(c(k))_{k\geq 0}$ and $(c^{*}(k))_{k\geq 0}$: 
equation~\eqref{equation_recursion_for_c*}
expresses $c(k)$ 
in terms of $c^{*}(0), \ldots, c^{*}(k-1)$, whereas the 
dual equation
expresses $c^{*}(k)$
in terms of $c(0), \ldots, c(k-1)$.
Together with the initial condition $c(0) = c^{*}(0) = 1$,
the two functions $C(x)$ and $C^{*}(x)$ are uniquely determined.


Now we prove the third identity.
Given a partition $\pi$ in $\NC^{d}_{dk+1}$,
let $j$ be the largest element in the block $B$ containing $1$.
Theorem~\ref{theorem_characterization} implies
that $j \equiv 1 \bmod d$, hence let $j = dr+1$.
Similarly, then $j^{\prime}$ is the smallest element in
block of the dual partition containing $n^{\prime}$.
Let $\sigma$ be the restriction $\pi$ to the interval $[1,j]$,
which has cardinality $dr+1$.
Let $\tau$ be the restriction of $\pi$ to the union $[j+1,n] \cup B$
and identify all elements of the block~$B$.
Note that $\tau$ is a partition on a set of size $d(k-r) + 1$.
Note that $\sigma$ has a given singleton block in the dual
and $\tau$ has also a given singleton block.
Enumerating the blocks in $\pi$ and the blocks in the dual partition
we have
\begin{align*}
|\pi| -1 & = |\sigma| -1 + |\tau| - 1 , &
|\pi^{\prime}| -1 & = |\sigma^{\prime}| -1 + |\tau^{\prime}| - 1 .
\end{align*}
Hence the weight of $\pi$ factors as
\begin{align*}
\wt(\pi) 
& =
\wt(\tau) \cdot \wt(\sigma) .
\end{align*}
Summing over all partitions $\pi$ in $\NC^{d}_{dk+1}$, multiplying with
$x^{k} = x^{r} \cdot x^{k-r}$ and summing over all
$k \geq 0$ yield the desired identity.
\end{proof}

\end{document}
